% ============================================================================
% CAPÍTULO 10: INVARIANTE DE BARGMANN-PANCHARATNAM Y FASE GEOMÉTRICA
% ============================================================================
\chapter{El Invariante de Bargmann-Pancharatnam: Fase Geométrica en $S^{D-1}$}


\textit{La fase que acumula un sistema cuántico al recorrer\\
un camino cerrado no depende del camino sino de la superficie\\
encerrada. Es la huella de la geometría sobre la física.} --- --- Berry, 1984

\section{Origen Físico: La Fase de Berry}

En 1984, Michael Berry descubrió que un sistema cuántico que evoluciona
\textit{adiabáticamente} a lo largo de un ciclo cerrado en el espacio de
parámetros $\mathcal{M}$ acumula una fase $\gamma$ que depende únicamente
de la geometría del camino, no de la velocidad de recorrido ni de la dinámica:

\begin{equation}
\gamma_B = i \oint_C \langle \psi(\mathbf{R}) | \nabla_{\mathbf{R}} | \psi(\mathbf{R}) \rangle \cdot d\mathbf{R}

\end{equation}

Esta fase no es un artefacto: es observable mediante interferometría y ha sido
medida en sistemas ópticos, espines nucleares, y superconductores.

\section{El Invariante de Bargmann-Pancharatnam}

La generalización discreta de la fase de Berry, debida a Bargmann (1964) y
Pancharatnam (1956), no requiere adiabaticidad ni evolución continua. Para una
secuencia discreta de estados $|\psi_0\rangle, |\psi_1\rangle, \ldots, |\psi_n\rangle$
con $|\psi_n\rangle = |\psi_0\rangle$ (ciclo cerrado):

\begin{equation}
\phi_{\text{BP}} = \arg\left(\langle\psi_0|\psi_1\rangle \langle\psi_1|\psi_2\rangle \cdots \langle\psi_{n-1}|\psi_n\rangle\right)

\end{equation}

\begin{theorem}[Invarianza del Invariante de Bargmann-Pancharatnam]
$\phi_{\text{BP}}$ es invariante bajo rephasing local: si se reemplaza
$|\psi_k\rangle \to e^{i\alpha_k}|\psi_k\rangle$ para escalares arbitrarios
$\alpha_k$, el producto de solapamientos satisface:
\begin{equation}
\prod_{k=0}^{n-1} \langle\psi_k| e^{-i\alpha_k} e^{i\alpha_{k+1}} |\psi_{k+1}\rangle = e^{i(\alpha_{n} - \alpha_0)} \prod_{k=0}^{n-1} \langle\psi_k|\psi_{k+1}\rangle
\end{equation}
Para un ciclo ($\alpha_n = \alpha_0$), los factores de fase se cancelan exactamente.
\end{theorem}

\section{Interpretación en $S^{D-1}$: Curvatura Acumulada}

En POLYDIM, los estados son puntos $\mathbf{x}_k \in S^{D-1}$ (vectores unitarios reales).
El solapamiento cuántico $\langle\psi_k|\psi_{k+1}\rangle$ corresponde al producto
interior $\langle \mathbf{x}_k, \mathbf{x}_{k+1} \rangle \in [-1, 1]$.

Para un ciclo de $n$ rotaciones de Rodrigues sobre $S^{D-1}$, el invariante de Bargmann-Pancharatnam es:

\begin{equation}
\phi_{\text{POLYDIM}} = \arg\left(\prod_{k=0}^{n-1} \langle \mathbf{x}_k, \mathbf{x}_{k+1} \rangle\right)
= \sum_{k=0}^{n-1} \text{sgn}(\langle \mathbf{x}_k, \mathbf{x}_{k+1}\rangle) \cdot \theta_k

\end{equation}

donde $\theta_k = \arccos(\langle \mathbf{x}_k, \mathbf{x}_{k+1}\rangle)$ es la distancia
geodésica entre estados consecutivos.

\begin{proposition}[Invariante BP como Detector de Holonomía]
Si $\phi_{\text{POLYDIM}} \ne 0$ para un ciclo cerrado de rotaciones de Rodrigues,
entonces las rotaciones no conmutan y el ciclo produce holonomía --- el vector
de estado final no coincide con el inicial incluso si el camino regresa al mismo punto.
\end{proposition}

\section{Relevancia para POLYDIM: Detección de Deriva Acumulada}

La holonomía en $S^{D-1}$ es un problema práctico en pipelines de rotación encadenada:
si se aplican $N$ rotaciones de Rodrigues sucesivas con el objetivo de implementar
una transformación compuesta, la deriva acumulada puede superar $\varepsilon_{\text{mach}}$
incluso cuando cada rotación individual tiene drift $\le \varepsilon_{\text{mach}}$.

El invariante de Bargmann-Pancharatnam permite cuantificar esta deriva sin comparar
directamente los vectores:

\begin{equation}
\text{Drift acumulado} \approx \frac{|\phi_{\text{BP}}|}{D} \cdot \text{Constante de normalización}
\end{equation}

\section{Conexión con el Invariante de Higham}

El Teorema 4.3 de Higham provee la cota de error \textit{local} por operación:
\begin{equation}
\text{Error local} \le (2D + 50)\varepsilon_{\text{mach}} \quad \text{(sin Neumaier)}
\end{equation}

Para $N$ operaciones encadenadas, la cota ingenuamente sería $N \times \text{Error local}$.
El invariante de Bargmann-Pancharatnam provee una cota \textit{geométrica} alternativa
que puede ser mucho más estricta cuando las rotaciones son casi conmutativas:

\begin{equation}
\text{Drift acumulado}_{\text{BP}} \le \frac{2}{\sqrt{D}} \cdot |\phi_{\text{BP}}| + N \cdot \varepsilon_{\text{mach}}
\end{equation}

\section{Aplicación Cuántica: Clifford Twirling como Anulación de BP}

El \textit{Clifford Twirling} (Randomized Compiling) de V764 tiene una interpretación
elegante en términos del invariante de Bargmann-Pancharatnam: al aleatorizar las
compuertas Clifford alrededor de cada compuerta T, se impide que las fases de BP
de errores sucesivos se acumulen coherentemente.

Sin twirling: $\phi_{\text{BP}}^{\text{error}} \propto N\varepsilon$ (acumulación lineal).
Con twirling: $\phi_{\text{BP}}^{\text{error}} \propto \sqrt{N}\varepsilon$ (acumulación aleatoria).

Esto es exactamente la transformación de ``ruido coherente a ruido incoherente'' descrita
en el Capítulo~\ref{ch:clifford_qpu}.
